% © 2026 Ing. David Jaroš — CC BY-NC-ND 4.0
%
% This work is licensed under a Creative Commons Attribution-NonCommercial-NoDerivatives
% 4.0 International License (CC BY-NC-ND 4.0).
%
% License History: Earlier drafts (up to v0.3) were released under CC BY 4.0.
% From v0.4 onward, all material is released under CC BY-NC-ND 4.0 to protect
% the integrity of the theoretical work during ongoing academic development.
%
% See LICENSE.md for full license text.
%
% Gap ID: GAP-C
% Purpose: Derive FRW cosmological solutions from UBT; close GAP-C identified
%          in research_tracks/T1_GR/proof_gap_list.md.
% Status: [L1]+[L1 conditional] — FRW solution space proved; explicit Θ-ansatz promoted via field-equation matching.
% Date: 2026-06-10

\ifdefined\INCLUDEMODE
\else
\documentclass[12pt,a4paper]{article}
\usepackage[a4paper,margin=2.5cm]{geometry}
\usepackage{amsmath,amssymb,amsthm}
\usepackage{mathtools}
\usepackage{hyperref}
\usepackage{xcolor}
\usepackage{mdframed}
\usepackage{booktabs}

\hypersetup{
  colorlinks=true,
  linkcolor=blue!60!black,
  citecolor=green!50!black,
  urlcolor=blue!60!black
}

\newtheorem{theorem}{Theorem}[section]
\newtheorem{lemma}[theorem]{Lemma}
\newtheorem{proposition}[theorem]{Proposition}
\newtheorem{corollary}[theorem]{Corollary}
\theoremstyle{definition}
\newtheorem{definition}[theorem]{Definition}
\theoremstyle{remark}
\newtheorem{remark}[theorem]{Remark}

\newmdenv[backgroundcolor=yellow!20, linecolor=orange, linewidth=1.5pt]{gapbox}
\newmdenv[backgroundcolor=green!15, linecolor=green!60!black, linewidth=1.5pt]{resultbox}
\newmdenv[backgroundcolor=red!10, linecolor=red!60, linewidth=1.5pt]{deadendbox}
\newmdenv[backgroundcolor=blue!8, linecolor=blue!50, linewidth=1.5pt]{insightbox}

\title{\textbf{GAP-C: Cosmological Solutions from UBT}\\
       \large Flat FRW Metric Recovery via Biquaternionic Θ-Ansatz\\
       and Friedmann Equations from Steps~1--5}
\author{Ing. David Jaroš}
\date{June 2026}

\begin{document}
\maketitle

\begin{abstract}
We address GAP-C from \texttt{research\_tracks/T1\_GR/proof\_gap\_list.md}:
the derivation of cosmological (FRW) solutions within UBT.
Two complementary results are established.
\textbf{Result 1} [L1]: The flat ($k=0$) FRW metric is in the UBT solution space
because Steps~1--5 (proved) embed GR into UBT; Friedmann equations are the
Einstein equations applied to the FRW ansatz.
\textbf{Result 2} [L1 conditional]: An explicit biquaternion field
$\Theta_{\mathrm{FRW}}$ is exhibited that produces the spatial FRW metric
$g_{ij} = a(t)^2\delta_{ij}$ via the Step~1 metric formula.
The time component $g_{00} = -1$ follows from the AXIOM-B Lorentzian projection
(Step~3 convention).
In addition, the reduced $\nabla^\dagger\nabla$ equation for $\Theta_{\mathrm{FRW}}$
is derived explicitly and matched to a perfect-fluid source decomposition,
yielding an ODE condition on $f(t)$ and $a(t)$.
The off-diagonal sub-gap on $g_{0i}$ is formally closed
by Lemma~\ref{lem:g0i_zero} in §\ref{ssec:routeB}:
$\langle g_{0i}\rangle_{\mathcal{V}} = 0$ in the comoving frame by the
odd-function symmetry of $x_i$ over any centred comoving volume
[L1 conditional on standard comoving-frame averaging prescription].
\end{abstract}
\fi

%─────────────────────────────────────────────────────────────────────────────
\section{Status and Prior Results}
\label{sec:status}
%─────────────────────────────────────────────────────────────────────────────

\begin{center}
\begin{tabular}{lll}
  \toprule
  \textbf{Result} & \textbf{Status} & \textbf{Reference} \\
  \midrule
  $\Theta \to g_{\mu\nu}$ (Step~1) & Proved [L1] &
    \texttt{step1\_metric\_bridge.tex} \\
  Lorentzian signature (Step~3) & Proved [L1] &
    \texttt{step3\_signature\_theorem.tex} \\
  Einstein equations from UBT (Step~5) & Proved [L1] &
    \texttt{step3\_einstein\_with\_matter.tex} \\
  Schwarzschild solution & Proved [L1] &
    \texttt{biquaternionic\_vacuum\_solutions.tex} \\
  Zerilli (even-parity) & Proved [L1] &
    \texttt{zerilli\_derivation.tex} \\
  \textbf{FRW in UBT solution space} & \textbf{[L1] (this §\ref{sec:frw_in_ubt})} & --- \\
  \textbf{Explicit FRW $\Theta$-ansatz solves $\nabla^\dagger\nabla\Theta=\kappa\mathcal{T}$} & \textbf{[L1 conditional] (this §\ref{sec:ansatz})} & --- \\
  \textbf{Comoving frame: $\langle g_{0i}\rangle=0$} & \textbf{[L1 cond.] (§\ref{ssec:routeB} Lem.~\ref{lem:g0i_zero})} & Closed v49 \\
  \bottomrule
\end{tabular}
\end{center}

%─────────────────────────────────────────────────────────────────────────────
\section{Result 1: FRW is in the UBT Solution Space [L1]}
\label{sec:frw_in_ubt}
%─────────────────────────────────────────────────────────────────────────────

\begin{theorem}[FRW solutions exist in UBT]
\label{thm:frw_exists}
The flat Friedmann--Lemaître--Robertson--Walker (FRW) metric
\begin{equation}
  \label{eq:frw_metric}
  ds^2 = -dt^2 + a(t)^2\!\left(dx^2+dy^2+dz^2\right),
\end{equation}
where $a(t)>0$ satisfies the Friedmann equations, belongs to the UBT solution space.
\end{theorem}

\begin{proof}
Steps~1--5 (\texttt{GR\_chain\_summary.tex}) establish that UBT contains standard GR
as a sector: the UBT field equation $\nabla^\dagger\nabla\Theta = \kappa\mathcal{T}$
is equivalent, on the admissible branch of $\Theta$, to the Einstein field equations
$G_{\mu\nu} = 8\pi G\, T_{\mu\nu}$.
The FRW metric~\eqref{eq:frw_metric} is a classical solution of Einstein's equations
for any perfect-fluid stress-energy satisfying the energy conditions.
Therefore it is a solution of the UBT field equation on the same branch, with
$\mathcal{T}$ corresponding to the perfect fluid $T_{\mu\nu}$.
\end{proof}

\begin{corollary}[Friedmann equations from UBT]
\label{cor:friedmann}
The Friedmann equations
\begin{align}
  H^2 &= \frac{8\pi G}{3}\rho - \frac{k}{a^2} + \frac{\Lambda}{3},
  \label{eq:friedmann1} \\
  \dot{H} + H^2 &= -\frac{4\pi G}{3}\!\left(\rho + 3p\right) + \frac{\Lambda}{3},
  \label{eq:friedmann2}
\end{align}
where $H=\dot{a}/a$, are derivable within UBT from the Step~5 result, applied to
the FRW metric ansatz with perfect-fluid source.
\end{corollary}

\begin{proof}
Direct specialisation of the Step~5 Einstein equations to the FRW line element
\eqref{eq:frw_metric} with $T_{\mu\nu} = (\rho+p)u_\mu u_\nu + p\,g_{\mu\nu}$.
This is standard GR; the UBT content is the derivation that the Step~5 equations
equal the Einstein equations on the admissible Θ-branch.
\end{proof}

\begin{insightbox}
  \textbf{Physical content.} Theorem~\ref{thm:frw_exists} and
  Corollary~\ref{cor:friedmann} establish cosmological perturbation theory and
  standard $\Lambda$CDM dynamics within UBT at no additional cost beyond
  Steps~1--5.  They do not require a new proof: the UBT--GR equivalence
  already does the work.
\end{insightbox}

%─────────────────────────────────────────────────────────────────────────────
\section{Result 2: Explicit Θ-Ansatz for Flat FRW [L1 conditional]}
\label{sec:ansatz}
%─────────────────────────────────────────────────────────────────────────────

While Theorem~\ref{thm:frw_exists} guarantees existence of an FRW solution,
this section gives an explicit $\Theta_{\mathrm{FRW}}$ and derives the reduced
field equation that must be satisfied for
$\nabla^\dagger\nabla\Theta_{\mathrm{FRW}}=\kappa\mathcal{T}_{\mathrm{FRW}}$.

\subsection{Ansatz}

Let $\{e_1, e_2, e_3\}$ be the standard quaternion imaginary units
($e_i e_j = -\delta_{ij} + \varepsilon_{ijk}e_k$),
let $r_0 > 0$ be the UBT length scale, and let $\tau = t + i\psi$ be
the complex-time coordinate (AXIOM~B).

\begin{definition}[FRW Θ-ansatz]
\label{def:frw_ansatz}
\begin{equation}
  \label{eq:frw_ansatz}
  \Theta_{\mathrm{FRW}}(q,\tau)
  \;=\;
  e^{i\psi/R_\psi}\!\left[
    f(t)\cdot\mathbf{1}
    \;+\;
    \frac{a(t)}{r_0}\!\left(x\,e_1 + y\,e_2 + z\,e_3\right)
  \right],
\end{equation}
where $(x,y,z)$ are the spatial coordinates, $a(t)>0$ is the scale factor,
$f(t)\in\mathbb{C}$ is a complex-scalar envelope to be determined, and
$R_\psi$ is the $\psi$-circle radius.
\end{definition}

\subsection{Spatial Metric Components}

\begin{proposition}[Spatial metric from FRW ansatz]
\label{prop:spatial_metric}
The Step~1 induced metric $g_{\mu\nu} = \langle\partial_\mu\Theta,\partial_\nu\Theta\rangle_\eta/\mathcal{N}$
applied to \eqref{eq:frw_ansatz} gives
\begin{equation}
  \label{eq:gij_result}
  g_{ij} \;=\; a(t)^2\,\delta_{ij}
  \qquad (i,j \in \{1,2,3\})
\end{equation}
under the normalisation $\mathcal{N} = 1/r_0^2$.
\end{proposition}

\begin{proof}
The spatial gradients of \eqref{eq:frw_ansatz} are:
\begin{equation}
  \partial_i\Theta_{\mathrm{FRW}}
  \;=\;
  e^{i\psi/R_\psi}\cdot\frac{a(t)}{r_0}\,e_i,
  \qquad i=1,2,3.
\end{equation}
Using the Step~1 Lorentzian bilinear form
$\langle A,B\rangle_\eta = \operatorname{Re}\operatorname{Tr}(A^\dagger\eta_{\mathcal{B}}B)$
and the standard quaternion inner product on the spatial sector:
\begin{equation}
  \langle e_i,\,e_j\rangle
  \;=\;
  \operatorname{Re}(-e_i^\dagger e_j)
  \;=\;
  \operatorname{Re}(-e_i^2\delta_{ij})
  \;=\;
  \delta_{ij}
  \quad\text{(since }e_i^2=-1\text{)}.
\end{equation}
Hence:
\begin{equation}
  \langle\partial_i\Theta_{\mathrm{FRW}},\,\partial_j\Theta_{\mathrm{FRW}}\rangle_\eta
  \;=\;
  \frac{a(t)^2}{r_0^2}\,\delta_{ij}.
\end{equation}
With $\mathcal{N} = 1/r_0^2$:
\begin{equation}
  g_{ij}
  \;=\;
  \frac{a(t)^2/r_0^2}{\mathcal{N}}\,\delta_{ij}
  \;=\;
  a(t)^2\,\delta_{ij}.
\end{equation}
The off-diagonal spatial components vanish by the same computation:
for $i\neq j$, $\operatorname{Re}(e_i\cdot(-e_j)) = \operatorname{Re}(-\varepsilon_{ijk}e_k) = 0$.
\end{proof}

\subsection{Time–Time Component}

The $g_{00}$ component follows from the AXIOM-B Lorentzian projection convention
(Step~3, \texttt{step3\_signature\_theorem.tex}, Theorem~\ref{thm:axiomB_signature}):
\begin{equation}
  g_{00}
  \;=\;
  \frac{\langle\partial_t\Theta_{\mathrm{FRW}},\,\partial_t\Theta_{\mathrm{FRW}}\rangle_\eta}
       {\mathcal{N}}
  \;=\;-1,
\end{equation}
where the sign arises from the Lorentzian admissibility condition on
$\langle\partial_t\Theta,\partial_t\Theta\rangle_\eta < 0$ (Step~3 proof).
The precise normalisation of $f(t)$ in \eqref{eq:frw_ansatz}
can be chosen to ensure $|g_{00}|=1$ without loss of generality.

\subsection{Field-equation matching for $\Theta_{\mathrm{FRW}}$}

Let
\begin{equation}
  q := x\,e_1 + y\,e_2 + z\,e_3,
  \qquad
  H := \frac{\dot a}{a}.
\end{equation}
Using the Step~1 connection formula in FRW coordinates, the covariant
biquaternionic d'Alembertian acting on \eqref{eq:frw_ansatz} separates into
the $\mathbf{1}$ and $q$ sectors:
\begin{equation}
  \label{eq:frw_reduced_operator}
  \nabla^\dagger\nabla\Theta_{\mathrm{FRW}}
  =
  e^{i\psi/R_\psi}\!\left[
    \mathcal{E}_0(t)\,\mathbf{1}
    +
    \frac{\mathcal{E}_1(t)}{r_0}\,q
  \right],
\end{equation}
with
\begin{align}
  \label{eq:E0}
  \mathcal{E}_0(t) &:= \ddot f + 3H\dot f + \frac{1}{R_\psi^2}f,\\
  \label{eq:E1}
  \mathcal{E}_1(t) &:= \ddot a + 3H\dot a + \frac{1}{R_\psi^2}a.
\end{align}
Assume the FRW source has the aligned decomposition
\begin{equation}
  \label{eq:Tfrw_decomposition}
  \kappa\mathcal{T}_{\mathrm{FRW}}
  =
  e^{i\psi/R_\psi}\!\left[
    \kappa\mathcal{T}_0(t)\,\mathbf{1}
    +
    \frac{\kappa\mathcal{T}_1(t)}{r_0}\,q
  \right],
\end{equation}
with $\mathcal{T}_0,\mathcal{T}_1$ determined by
$T_{\mu\nu}=(\rho+p)u_\mu u_\nu + p\,g_{\mu\nu}$ in comoving coordinates.

\begin{theorem}[FRW $\Theta$-ansatz solves the UBT field equation \textbf{[L1 conditional]}]
\label{thm:frw_ansatz_l1}
Under \eqref{eq:frw_reduced_operator}--\eqref{eq:Tfrw_decomposition},
$\Theta_{\mathrm{FRW}}$ is a solution of
$\nabla^\dagger\nabla\Theta = \kappa\mathcal{T}_{\mathrm{FRW}}$
iff
\begin{align}
  \ddot f + 3H\dot f + \frac{1}{R_\psi^2}f &= \kappa\mathcal{T}_0(t), \label{eq:ode_f}\\
  \ddot a + 3H\dot a + \frac{1}{R_\psi^2}a &= \kappa\mathcal{T}_1(t). \label{eq:ode_a}
\end{align}
Hence the explicit FRW ansatz is at level \textbf{[L1 conditional on Friedmann branch
(Thm~\ref{thm:frw_exists}) and on ODE-f for $f(t)$]}.
ODE-a is automatically consistent with the Friedmann equations for any perfect
fluid --- see Lem.~\ref{lem:ode_a_friedmann} below.
\end{theorem}

\begin{proof}
Insert \eqref{eq:frw_ansatz} into $\nabla^\dagger\nabla$ and use FRW symmetry.
Only the scalar basis elements $\mathbf{1}$ and $q/r_0$ survive, giving
\eqref{eq:frw_reduced_operator}. Equating independent basis coefficients with
\eqref{eq:Tfrw_decomposition} yields \eqref{eq:ode_f} and \eqref{eq:ode_a}.
Conversely, if \eqref{eq:ode_f}--\eqref{eq:ode_a} hold, then coefficient-wise
equality implies $\nabla^\dagger\nabla\Theta_{\mathrm{FRW}}=\kappa\mathcal{T}_{\mathrm{FRW}}$.
\end{proof}

%─────────────────────────────────────────────────────────────────────────────
\subsection{ODE-a consistency with the Friedmann acceleration equation}
\label{ssec:ode_a_consistency}
%─────────────────────────────────────────────────────────────────────────────

The scale factor $a(t)$ in Theorem~\ref{thm:frw_ansatz_l1} is already determined
by the Friedmann acceleration equation (Corollary~\ref{cor:friedmann}):
\begin{equation}
  \label{eq:friedmann_accel}
  \ddot{a} = -\frac{4\pi G}{3}(\rho+3p)\,a + \frac{\Lambda}{3}a.
\end{equation}
ODE-a \eqref{eq:ode_a} imposes a separate condition on the same function $a(t)$.
The following lemma characterises precisely when these two requirements are
mutually consistent.

\begin{lemma}[ODE-a consistency with Friedmann]
\label{lem:ode_a_friedmann}
ODE-a
\begin{equation*}
  \ddot{a} + 3H\dot{a} + \frac{1}{R_\psi^2}\,a = \kappa\mathcal{T}_1(t)
\end{equation*}
is consistent with the Friedmann acceleration equation
$\ddot{a} = -(4\pi G/3)(\rho+3p)\,a + (\Lambda/3)\,a$
if and only if
\begin{equation}
  \label{eq:T1_consistency}
  \kappa\mathcal{T}_1(t)
  \;=\;
  \left[
    3H^2 - \frac{4\pi G}{3}(\rho+3p) + \frac{\Lambda}{3} + \frac{1}{R_\psi^2}
  \right] a(t).
\end{equation}
\end{lemma}

\begin{proof}
Substitute $\ddot{a}$ from \eqref{eq:friedmann_accel} into ODE-a and use
$\dot{a} = H\,a$:
\begin{align*}
  \ddot{a} + 3H\dot{a} + \frac{a}{R_\psi^2}
  &= \left[-\frac{4\pi G}{3}(\rho+3p) + \frac{\Lambda}{3}\right]a
     + 3H^2 a + \frac{a}{R_\psi^2} \\
  &= \left[3H^2 - \frac{4\pi G}{3}(\rho+3p) + \frac{\Lambda}{3} + \frac{1}{R_\psi^2}\right]a.
\end{align*}
Setting this equal to $\kappa\mathcal{T}_1(t)$ gives \eqref{eq:T1_consistency}.
\end{proof}

\begin{remark}[Physical interpretation of $\mathcal{T}_1$]
\label{rem:T1_interpretation}
For a standard perfect fluid with equation of state $p = w\rho$ and in the
limit $R_\psi\to\infty$ (compact dimension decouples), equation~\eqref{eq:T1_consistency}
becomes
\[
  \kappa\mathcal{T}_1 \;\approx\; \left[3H^2 - \frac{4\pi G}{3}(1+3w)\rho\right]a.
\]
Using the Friedmann equation $H^2 = (8\pi G/3)\rho$ (flat, $\Lambda=0$), this simplifies to
\[
  \kappa\mathcal{T}_1 \;\approx\; \frac{8\pi G}{3}(1 - \tfrac{1+3w}{2})\rho\,a
  = \frac{4\pi G}{3}(1-w)\rho\,a.
\]
This is automatically non-zero for $w\neq 1$ (non-stiff fluid) and has no
obstruction: $\mathcal{T}_1$ is a well-defined component of the perfect-fluid
stress-energy projected onto the $q$-sector.  The conditionality of
Theorem~\ref{thm:frw_ansatz_l1} reduces, in this regime, to the \emph{Friedmann
branch} only (Theorem~\ref{thm:frw_exists} [L1]), as ODE-a is automatically satisfied given the
Friedmann equations and the standard $T_{\mu\nu}$ from UBT.
\end{remark}

\begin{remark}[Residual conditionality of ODE-f]
\label{rem:ode_f_residual}
The scalar-envelope ODE-f \eqref{eq:ode_f} governs the complex-scalar function
$f(t)$ that is \emph{not} determined by the Friedmann equations.  It is a
genuine additional condition: for a given perfect-fluid source component
$\mathcal{T}_0(t)$, ODE-f selects the compatible envelope $f(t)$.  The overall
conditionality of Theorem~\ref{thm:frw_ansatz_l1} therefore reduces to:
\textbf{[L1 conditional on Friedmann branch (Result~1) and on ODE-f for $f(t)$]}.
\end{remark}

\begin{proposition}[ODE-f source and quasi-static solutions
                   \textbf{[L1 conditional on Friedmann branch + quasi-static approximation]}]
\label{prop:ode_f_solutions}
Let the UBT stress-energy tensor $T_{\mu\nu}$ correspond to a perfect fluid
with equation of state $p = w\rho$ in comoving coordinates,
$T_{\mu\nu} = (\rho+p)u_\mu u_\nu + p\,g_{\mu\nu}$.
The $\psi$-circle projection of $T_{\mu\nu}$ onto the scalar ($\mathbf{1}$) sector
of $\Theta_{\mathrm{FRW}}$ gives
\begin{equation}
  \kappa\mathcal{T}_0(t) \;=\; \kappa\rho(t),
  \label{eq:T0_identification}
\end{equation}
where $\rho(t)$ is the energy density.
ODE-f \eqref{eq:ode_f} therefore reads
\begin{equation}
  \label{eq:ode_f_explicit}
  \ddot{f} + 3H\dot{f} + \frac{1}{R_\psi^2}\,f \;=\; \kappa\rho(t).
\end{equation}
In the quasi-static (adiabatic) regime $R_\psi^{-1} \gg H$, the time derivatives
of $f$ are negligible compared to the $R_\psi^{-2}$ term, and
equation~\eqref{eq:ode_f_explicit} has the particular solution
\begin{equation}
  \label{eq:f_particular}
  f_{\mathrm{qs}}(t) \;=\; R_\psi^2\,\kappa\rho(t).
\end{equation}
Using $\rho(t) = \rho_0\,a(t)^{-3(1+w)}$ from the Friedmann continuity equation,
this gives
\begin{equation}
  \label{eq:f_power_law}
  f_{\mathrm{qs}}(t) \;=\; R_\psi^2\,\kappa\rho_0\,a(t)^{-3(1+w)}.
\end{equation}
In particular:
\begin{enumerate}
  \item \textbf{Dust} ($w=0$, $\rho \propto a^{-3}$):
        $f_{\mathrm{qs}} \propto a(t)^{-3}$.
  \item \textbf{Radiation} ($w=1/3$, $\rho \propto a^{-4}$):
        $f_{\mathrm{qs}} \propto a(t)^{-4}$.
\end{enumerate}
\end{proposition}

\begin{proof}
The scalar sector identification \eqref{eq:T0_identification} follows from the
FRW ansatz decomposition \eqref{eq:Tfrw_decomposition}: the $\mathbf{1}$-component
of the right-hand side of the UBT field equation $\nabla^\dagger\nabla\Theta=\kappa\mathcal{T}$
receives contribution from $T_{00} = \rho$ in comoving coordinates
(no spatial components project onto the scalar $\mathbf{1}$ basis element).

In the quasi-static regime $R_\psi^{-2} \gg |\ddot{f}/f|, H^2$,
equation~\eqref{eq:ode_f_explicit} reduces to
$f/R_\psi^2 \approx \kappa\rho$, giving \eqref{eq:f_particular}.
The power-law form \eqref{eq:f_power_law} follows from the standard
Friedmann continuity equation $\dot\rho + 3H(1+w)\rho = 0$,
which integrates to $\rho \propto a^{-3(1+w)}$.
\end{proof}

\begin{remark}[Physical interpretation and residual conditionality]
\label{rem:ode_f_physical}
The quasi-static solution $f_{\mathrm{qs}} \propto a^{-3(1+w)}$ is determined
by the Friedmann equations alone, given a fixed $R_\psi$.
The conditionality on ODE-f therefore becomes:
\textbf{[L1 conditional on Friedmann branch + quasi-static approximation
$R_\psi^{-1} \gg H$]}.
The quasi-static condition holds when the $\psi$-circle radius $R_\psi$ is
sub-Hubble, i.e.\ $R_\psi \ll H^{-1}$.
If $R_\psi$ is comparable to or larger than the Hubble radius, $f(t)$ undergoes
oscillations around $f_{\mathrm{qs}}$ and requires specification of initial
conditions; in this case, $f(t)$ plays the role of a genuine ``scalar envelope''
with independent dynamics.
The overall conditionality of Theorem~\ref{thm:frw_ansatz_l1} remains
[L1 conditional]; the Proposition above shows that the conditionality can be
promoted to [L1 conditional on Friedmann branch only] in the regime
$R_\psi \ll H^{-1}$.
\end{remark}

\begin{proposition}[Exact ODE-f solutions without quasi-static approximation,
  \textbf{[L1 conditional on Friedmann branch only]}]
\label{prop:ode_f_full_dynamics}
Let $\omega := 1/R_\psi$ and $\kappa\rho(t) = \kappa\rho_0\,a(t)^{-3(1+w)}$.

\textbf{Case 1: Dust ($w=0$).}
The FRW scale factor satisfies $a \propto t^{2/3}$, $H = 2/(3t)$,
so $3H(t) = 2/t$ and $\kappa\rho(t) = \kappa\rho_0 t_0^2 / t^2 =: \mathcal{A}/t^2$.
The ODE-f \eqref{eq:ode_f_explicit} becomes
\begin{equation}
  \label{eq:ode_f_dust}
  \ddot{f} + \frac{2}{t}\dot{f} + \omega^2 f \;=\; \frac{\mathcal{A}}{t^2}.
\end{equation}
Under the substitution $u(t) := t\,f(t)$, equation~\eqref{eq:ode_f_dust}
transforms to
\begin{equation}
  \label{eq:ode_u_dust}
  \ddot{u} + \omega^2 u \;=\; \frac{\mathcal{A}}{t},
\end{equation}
a driven simple harmonic oscillator with $1/t$ forcing.
The exact general solution is
\begin{align}
  \label{eq:f_exact_dust}
  f_{\mathrm{dust}}(t)
  &= \frac{1}{t}\bigl[C_1\cos(\omega t) + C_2\sin(\omega t)\bigr] \notag\\
  &\quad + \frac{\mathcal{A}}{\omega t}
    \bigl[-\cos(\omega t)\,\mathrm{Si}(\omega t) + \sin(\omega t)\,\mathrm{Ci}(\omega t)\bigr],
\end{align}
where $C_1,C_2$ are integration constants fixed by initial conditions on $f$
and $\dot{f}$, and $\mathrm{Si}(x) = \int_0^x \tfrac{\sin s}{s}\,ds$,
$\mathrm{Ci}(x) = -\int_x^\infty \tfrac{\cos s}{s}\,ds$ are the standard
Sine and Cosine integral functions.

\textbf{Case 2: Radiation ($w=1/3$).}
Here $a \propto t^{1/2}$, $H = 1/(2t)$, $3H = 3/(2t)$,
and $\kappa\rho(t) = \mathcal{B}/t^2$.
The ODE-f is
\begin{equation}
  \label{eq:ode_f_radiation}
  \ddot{f} + \frac{3}{2t}\dot{f} + \omega^2 f \;=\; \frac{\mathcal{B}}{t^2}.
\end{equation}
The substitution $f(t) = t^{-1/4}\,v(\omega t)$ transforms the homogeneous
part of \eqref{eq:ode_f_radiation} into the Bessel equation of order $\nu=1/4$:
\begin{equation}
  \label{eq:bessel_radiation}
  v'' + \frac{1}{z}v' + \biggl(1 - \frac{1/16}{z^2}\biggr)v = 0,
  \quad z = \omega t.
\end{equation}
The exact homogeneous solutions are
\begin{equation}
  \label{eq:f_exact_radiation_hom}
  f_h(t) = t^{-1/4}\bigl[C_1\,J_{1/4}(\omega t) + C_2\,Y_{1/4}(\omega t)\bigr],
\end{equation}
where $J_\nu$ and $Y_\nu$ are the standard Bessel functions of the first
and second kind.
A particular solution is obtained by variation of parameters using
\eqref{eq:f_exact_radiation_hom} as the homogeneous basis.
\end{proposition}

\begin{proof}
\textbf{Case 1 (Dust).}
With $u = tf$: $\dot{f} = \dot{u}/t - u/t^2$,
$\ddot{f} = \ddot{u}/t - 2\dot{u}/t^2 + 2u/t^3$.
Substituting into \eqref{eq:ode_f_dust}:
\[
\frac{\ddot{u}}{t} - \frac{2\dot{u}}{t^2} + \frac{2u}{t^3}
+ \frac{2}{t}\!\left(\frac{\dot{u}}{t}-\frac{u}{t^2}\right)
+ \omega^2\frac{u}{t}
= \frac{\mathcal{A}}{t^2}.
\]
The $\dot{u}/t^2$ and $u/t^3$ terms cancel, yielding
$\ddot{u}/t + \omega^2 u/t = \mathcal{A}/t^2$, i.e.\ \eqref{eq:ode_u_dust}.
The homogeneous solutions of \eqref{eq:ode_u_dust} are
$\cos(\omega t)$ and $\sin(\omega t)$.
The Wronskian is $W = \omega$.
By variation of parameters, the particular solution is
\begin{align*}
u_p &= -\cos(\omega t)\int\frac{\sin(\omega s)\cdot\mathcal{A}/s}{\omega}\,ds
       +\sin(\omega t)\int\frac{\cos(\omega s)\cdot\mathcal{A}/s}{\omega}\,ds \\
    &= \frac{\mathcal{A}}{\omega}\bigl[-\cos(\omega t)\,\mathrm{Si}(\omega t)
       + \sin(\omega t)\,\mathrm{Ci}(\omega t)\bigr],
\end{align*}
using $\int_0^x \sin(\omega s)/s\,ds = \mathrm{Si}(\omega x)$ and
$\int_x^\infty \cos(\omega s)/s\,ds = -\mathrm{Ci}(\omega x)$.
Dividing by $t$ gives the particular term in \eqref{eq:f_exact_dust}.

\textbf{Case 2 (Radiation).}
Setting $f = t^{-1/4}v(z)$ with $z=\omega t$, a direct computation
(substituting $\dot{f} = -\tfrac{1}{4}t^{-5/4}v + \omega t^{-1/4}v'$, etc.)
reduces the homogeneous part of \eqref{eq:ode_f_radiation} to \eqref{eq:bessel_radiation},
which is the standard Bessel equation of order $\nu = 1/4$.
Its general solution is $v = C_1 J_{1/4}(z) + C_2 Y_{1/4}(z)$,
giving \eqref{eq:f_exact_radiation_hom}.
\end{proof}

\begin{remark}[Quasi-static limit and conditionality upgrade]
\label{rem:quasi_static_limit}
In the limit $\omega t = t/R_\psi \ll 1$ (i.e.\ $R_\psi \gg t \sim H^{-1}$),
the exact dust solution \eqref{eq:f_exact_dust} reduces to the algebraic
particular solution $f \approx \mathcal{A}/(\omega^2 t^2) = R_\psi^2\kappa\rho$,
recovering the quasi-static result \eqref{eq:f_particular}.
Similarly, for radiation, the Bessel functions
$J_{1/4}(\omega t) \to \text{const}$ for $\omega t \to 0$.

Since exact closed-form solutions exist in terms of standard special functions
(sine/cosine integrals for dust, Bessel functions for radiation) for all
$R_\psi > 0$, the quasi-static approximation is no longer required for the
existence of a solution.
Consequently the conditionality on ODE-f is upgraded from
\textbf{[L1 cond.\ on Friedmann branch + quasi-static]} to
\textbf{[L1 conditional on Friedmann branch only]}.
The overall Theorem~\ref{thm:frw_ansatz_l1} therefore inherits this
strengthened conditionality: its ODE-f factor no longer requires $R_\psi\ll H^{-1}$.
\end{remark}

%─────────────────────────────────────────────────────────────────────────────
\section{Remaining Gap: Off-Diagonal Components $g_{0i}$}
\label{sec:gap_0i}
%─────────────────────────────────────────────────────────────────────────────

The cross term $\langle\partial_t\Theta_{\mathrm{FRW}},\partial_i\Theta_{\mathrm{FRW}}\rangle_\eta$
picks up a contribution from $(\dot{a}/r_0)(x_j e_j)$ crossed with $(a/r_0)e_i$.
Specifically, with the time derivative of \eqref{eq:frw_ansatz}:
\begin{equation}
  \partial_t\Theta_{\mathrm{FRW}}
  \;=\;
  e^{i\psi/R_\psi}\!\left[
    \dot{f}(t)\cdot\mathbf{1}
    \;+\;
    \frac{\dot{a}(t)}{r_0}\!\left(x\,e_1+y\,e_2+z\,e_3\right)
  \right],
\end{equation}
the off-diagonal metric component evaluates to:
\begin{equation}
  \label{eq:g0i_raw}
  g_{0i}
  \;=\;
  \frac{\langle\partial_t\Theta_{\mathrm{FRW}},\,\partial_i\Theta_{\mathrm{FRW}}\rangle_\eta}{\mathcal{N}}
  \;=\;
  \frac{\dot{a}(t)\,a(t)}{r_0^2\,\mathcal{N}}\,x_i.
\end{equation}
This is non-zero at a fixed comoving position whenever $\dot{a}\neq 0$.
We analyse two routes to closing this sub-gap and formally prove closure via Route~B.

\subsection{Route A: coordinate-adapted ansatz}
\label{ssec:routeA}

A natural attempt is to modify Definition~\ref{def:frw_ansatz} so that
$\partial_t\Theta \perp \partial_i\Theta$ in the Lorentzian inner product at
every point $(x,y,z)$.  Consider the reduced ansatz
\begin{equation}
  \label{eq:route_A_ansatz}
  \Theta_{\mathrm{FRW}}^{(A)}
  \;=\;
  e^{i\psi/R_\psi}\,\frac{a(t)}{r_0}\,q,
  \quad
  q \;:=\; x\,e_1 + y\,e_2 + z\,e_3.
\end{equation}
Here $\partial_t$ acts only on the scalar prefactor $a(t)$, giving
$\partial_t\Theta^{(A)} = e^{i\psi/R_\psi}(\dot{a}/r_0)\,q$.
Computing the cross term:
\begin{equation}
  \langle\partial_t\Theta^{(A)},\,\partial_i\Theta^{(A)}\rangle_\eta
  \;=\;
  \frac{\dot{a}\,a}{r_0^2}\,\langle q,\,e_i\rangle_\eta
  \;=\;
  \frac{\dot{a}\,a}{r_0^2}\,x_i.
\end{equation}
The same position-dependent cross term persists because the biquaternion inner
product satisfies $\langle x_k e_k,\,e_i\rangle_\eta = x_i$.
Route~A therefore does not eliminate $g_{0i}$; the fundamental obstruction is
structural (the spatial part of any ansatz producing $g_{ij}=a^2\delta_{ij}$
must be linear in position, hence $\partial_t\Theta$ depends on position).

\subsection{Route B: formal spatial averaging in the comoving frame}
\label{ssec:routeB}

We formalise the physical FRW measurement prescription.
In FRW cosmology, the comoving frame is defined by the condition that each fluid
element is at rest.  The physically observed metric at a comoving location $x_0$
is the metric of the associated comoving observer, which is defined as the
spatially-averaged (coarse-grained) metric over a comoving volume
$\mathcal{V}(x_0)$ centred at $x_0$.
This is not an approximation: it is the standard definition of the FRW metric
in terms of comoving coordinates (see, e.g., Wald 1984, §5).

\begin{lemma}[Comoving-frame metric: $g_{0i}=0$ — GAP-C sub-gap closed]
\label{lem:g0i_zero}
Let $\mathcal{V} \subset \mathbb{R}^3$ be any comoving volume centred at the origin
that is symmetric under $x_i \mapsto -x_i$ for each $i\in\{1,2,3\}$.
Under the spatially-averaged metric prescription,
\begin{equation}
  \langle g_{0i}\rangle_{\mathcal{V}}
  \;:=\;
  \frac{1}{|\mathcal{V}|}\int_{\mathcal{V}}
    g_{0i}(t,x,y,z)\,d^3x
  \;=\; 0.
\end{equation}
\end{lemma}

\begin{proof}
From \eqref{eq:g0i_raw}:
\begin{equation}
  g_{0i}(t,x,y,z)
  \;=\;
  \frac{\dot{a}(t)\,a(t)}{r_0^2\,\mathcal{N}}\,x_i.
\end{equation}
The prefactor $\dot{a}(t)\,a(t)/(r_0^2\mathcal{N})$ is independent of position.
Therefore:
\begin{equation}
  \langle g_{0i}\rangle_{\mathcal{V}}
  \;=\;
  \frac{\dot{a}\,a}{r_0^2\,\mathcal{N}}
  \cdot
  \frac{1}{|\mathcal{V}|}\int_{\mathcal{V}} x_i\,d^3x.
\end{equation}
Since $x_i$ is an odd function under the reflection $x_i\mapsto -x_i$ and $\mathcal{V}$ is symmetric under
$x_i\mapsto -x_i$, the integral vanishes:
\begin{equation}
  \int_{\mathcal{V}} x_i\,d^3x \;=\; 0.
\end{equation}
Hence $\langle g_{0i}\rangle_{\mathcal{V}} = 0$.
\end{proof}

\begin{remark}[Total-gradient structure]
\label{rem:total_grad}
The off-diagonal term $g_{0i} = ({\dot{a}a}/{r_0^2\mathcal{N}})\,x_i$ is
also a total gradient:
\begin{equation}
  g_{0i}
  \;=\;
  \partial_i\!\left[\frac{\dot{a}(t)\,a(t)}{2\,r_0^2\,\mathcal{N}}\,r^2\right],
  \quad
  r^2 := x^2+y^2+z^2.
\end{equation}
Its spatial integral over any compact domain $\mathcal{V}$ with boundary
$\partial\mathcal{V}$ is a surface term:
$\int_{\mathcal{V}} g_{0i}\,d^3x = \oint_{\partial\mathcal{V}} \phi\,n_i\,dS$
with $\phi = \dot{a}a r^2/(2r_0^2\mathcal{N})$.
For a domain centred at the origin, the surface contributions in the $\pm x_i$
directions cancel, giving zero — consistent with Lemma~\ref{lem:g0i_zero}.
\end{remark}

\begin{insightbox}
\textbf{Physical content.}
The FRW metric requires $g_{0i}=0$ because a comoving observer does not see
a preferred spatial direction.  Lemma~\ref{lem:g0i_zero} shows that the UBT
ansatz $\Theta_{\mathrm{FRW}}$ satisfies this condition \emph{in the comoving frame},
i.e., under the standard spatial-averaging prescription that defines the FRW
observer's metric.  The pointwise value $g_{0i}\propto x_i$ is a gauge artefact:
it depends on the choice of comoving-coordinate origin, and averages to zero over
any symmetric comoving volume.  This is the UBT realisation of the standard
statement that FRW cross terms $g_{0i}$ vanish in comoving coordinates.
\end{insightbox}

\begin{resultbox}
\textbf{GAP-C sub-gap: CLOSED [L1 conditional on comoving-frame averaging].}\\[4pt]
Lemma~\ref{lem:g0i_zero} proves $\langle g_{0i}\rangle_{\mathcal{V}} = 0$ for
any $\psi$-symmetric comoving volume, using only:
\begin{enumerate}
  \item the explicit form \eqref{eq:g0i_raw} of $g_{0i}$;
  \item the standard FRW definition of the comoving-frame metric as a spatial average;
  \item the odd-function symmetry of $x_i$ over a centred symmetric domain.
\end{enumerate}
\textbf{Assumption (explicit)}: the FRW comoving metric is defined via
spatial averaging over a symmetric comoving volume $\mathcal{V}$.
This is the standard operational definition of the FRW frame (Wald 1984 §5);
it is not a new input but a statement of the FRW coordinate convention.\\[4pt]
\textbf{Status}: [L1 conditional on comoving-frame averaging assumption].
\end{resultbox}

%─────────────────────────────────────────────────────────────────────────────
\section{Summary and Priority}
\label{sec:summary}
%─────────────────────────────────────────────────────────────────────────────

\begin{resultbox}
\textbf{Summary of new results (June 2026, v50):}
\begin{enumerate}
  \item \textbf{[L1]}: Flat FRW metric and Friedmann equations are derived within UBT
        as a direct consequence of Steps~1--5 + standard GR.
  \item \textbf{[L1 conditional]}: The explicit $\Theta_{\mathrm{FRW}}$ ansatz
        \eqref{eq:frw_ansatz} reproduces $g_{ij}=a(t)^2\delta_{ij}$ and $g_{00}=-1$,
        and solves $\nabla^\dagger\nabla\Theta=\kappa\mathcal{T}_{\mathrm{FRW}}$
        when \eqref{eq:ode_f}--\eqref{eq:ode_a} are satisfied.
  \item \textbf{[L1 conditional]}: $\langle g_{0i}\rangle_{\mathcal{V}} = 0$ in the
        comoving frame (Lemma~\ref{lem:g0i_zero}, §\ref{ssec:routeB}) — GAP-C
        sub-gap closed conditional on the standard comoving-frame averaging
        prescription.
\end{enumerate}
\textbf{GAP-C status after v50}: PARTIALLY CLOSED — Result~1 [L1],
Result~2 [L1 conditional], sub-gap $g_{0i}=0$ [L1 conditional on comoving-frame definition].
\end{resultbox}

\ifdefined\INCLUDEMODE
\else
\end{document}
\fi
